% Fichier engendré par tools/generate-tex.py à partir de 07-algorithmique-et-programmation.
% Ne pas éditer à la main : tout le texte provient des commentaires du fichier Lean.
\documentclass[11pt,a4paper]{article}

% Compile aussi bien avec pdflatex qu'avec xelatex ou tectonic.
\usepackage{iftex}
\ifPDFTeX
  \usepackage[T1]{fontenc}
  \usepackage[utf8]{inputenc}
\else
  \usepackage{fontspec}
\fi
\usepackage[french]{babel}
\usepackage{amsmath,amssymb,amsthm}
\usepackage[margin=2.5cm]{geometry}
\usepackage[hidelinks]{hyperref}

% Les figures sont dessinées en TikZ : elles reprennent ainsi les
% polices du document. Voir la skill `illustrate-theorem`.
\usepackage{tikz}
\usetikzlibrary{calc}

\theoremstyle{definition}
\newtheorem{definition}{Définition}
\theoremstyle{plain}
\newtheorem{theoreme}{Théorème}
\newtheorem{lemme}[theoreme]{Lemme}

% Renvoi à la déclaration Lean, une ligne après la démonstration, aligné à droite.
\newcommand{\source}[2]{\par\smallskip\noindent\hfill{\footnotesize\href{#1}{\texttt{#2}}}\par\smallskip}

\title{Algorithmique et programmation}
\date{}

\begin{document}
\maketitle
% Les identifiants Lean en machine à écrire ne se coupent pas : on tolère des
% espacements plus lâches plutôt que des lignes qui débordent.
\sloppy

\section{AlgorithmiqueEtProgrammation}

\noindent
Collège \textemdash{} section \og{} Algorithmique et programmation \fg{}.
Un programme est ici une suite d'instructions agissant sur la valeur d'une variable.
Énoncés et démonstrations en français : voir AlgorithmiqueEtProgrammation.tex.

\subsection{Définitions}

\noindent
Toutes les définitions du fichier sont posées ici, avant les démonstrations.

\begin{definition}
Une \emph{instruction} transforme la valeur de la variable : c'est une
application de $\mathbb{Z}$ dans $\mathbb{Z}$.
\end{definition}
\source{https://github.com/Commutator-IO/learn-lean/blob/main/courses/01-college/07-algorithmique-et-programmation/AlgorithmiqueEtProgrammation.lean\#L15}{AlgorithmiqueEtProgrammation.lean\#L15}

\begin{definition}
\emph{Exécuter} une suite d'instructions, c'est les appliquer de la première à la
dernière :
\[ \mathrm{ex}\bigl([f_1, \dots, f_n]\bigr)(x) = f_n\bigl(\cdots f_1(x) \cdots\bigr) . \]
\end{definition}
\source{https://github.com/Commutator-IO/learn-lean/blob/main/courses/01-college/07-algorithmique-et-programmation/AlgorithmiqueEtProgrammation.lean\#L18}{AlgorithmiqueEtProgrammation.lean\#L18}

\begin{definition}
\emph{Répéter} $n$ fois une instruction, c'est l'itérer : $r_f(n) = f^{\,n}$, avec
$f^{\,0} = \mathrm{id}$ et $f^{\,n+1} = f \circ f^{\,n}$.
\end{definition}
\source{https://github.com/Commutator-IO/learn-lean/blob/main/courses/01-college/07-algorithmique-et-programmation/AlgorithmiqueEtProgrammation.lean\#L21}{AlgorithmiqueEtProgrammation.lean\#L21}

\begin{definition}
La \emph{boucle} « retrancher $b$ tant que c'est possible » :
\[ \mathrm{stq}(b, x) =
   \begin{cases}
     \mathrm{stq}(b, x - b) & \text{si } 0 < b \text{ et } b \le x, \\
     x & \text{sinon.}
   \end{cases} \]
\end{definition}
\source{https://github.com/Commutator-IO/learn-lean/blob/main/courses/01-college/07-algorithmique-et-programmation/AlgorithmiqueEtProgrammation.lean\#L24}{AlgorithmiqueEtProgrammation.lean\#L24}

\begin{definition}
Le programme $A$ : « ajouter 3, doubler, retrancher 6 ».
\end{definition}
\source{https://github.com/Commutator-IO/learn-lean/blob/main/courses/01-college/07-algorithmique-et-programmation/AlgorithmiqueEtProgrammation.lean\#L29}{AlgorithmiqueEtProgrammation.lean\#L29}

\begin{definition}
Le programme $B$ : « doubler ».
\end{definition}
\source{https://github.com/Commutator-IO/learn-lean/blob/main/courses/01-college/07-algorithmique-et-programmation/AlgorithmiqueEtProgrammation.lean\#L32}{AlgorithmiqueEtProgrammation.lean\#L32}

\subsection{Une suite d'affectations compose ses instructions dans l'ordre d'écriture}

\begin{theoreme}
Une suite d'affectations compose ses instructions dans l'ordre d'écriture :
\[ \mathrm{ex}\bigl([x \mapsto x+3,\ x \mapsto 2x]\bigr)(x) = (x + 3) \times 2 . \]
\end{theoreme}

\begin{proof}
On applique les instructions l'une après l'autre : la première donne $x + 3$, la seconde
double ce résultat. La valeur finale se lit ainsi de proche en proche.
\end{proof}
\source{https://github.com/Commutator-IO/learn-lean/blob/main/courses/01-college/07-algorithmique-et-programmation/AlgorithmiqueEtProgrammation.lean\#L38}{AlgorithmiqueEtProgrammation.lean\#L38}

\begin{theoreme}
Ajouter une instruction à la fin d'un programme, c'est l'appliquer au résultat :
\[ \mathrm{ex}(p \mathbin{+\!\!+} [f])(x) = f\bigl(\mathrm{ex}(p)(x)\bigr) . \]
\end{theoreme}

\begin{proof}
L'exécution parcourt la liste de gauche à droite : arrivée au bout de $p$, elle tient la
valeur $\mathrm{ex}(p)(x)$, à laquelle elle applique la dernière instruction.
\end{proof}
\source{https://github.com/Commutator-IO/learn-lean/blob/main/courses/01-college/07-algorithmique-et-programmation/AlgorithmiqueEtProgrammation.lean\#L43}{AlgorithmiqueEtProgrammation.lean\#L43}

\subsection{Boucle bornée : répéter \texttt{n} fois une instruction}

\begin{theoreme}
Boucle bornée : répéter $n$ fois l'ajout de $a$ ajoute $n \times a$ :
\[ \bigl(x \mapsto x + a\bigr)^{\,n}(x) = x + na . \]
\end{theoreme}

\begin{proof}
Par récurrence sur $n$. Pour $n = 0$, la boucle ne fait rien et $x + 0 \times a = x$.
Si le résultat vaut pour $n$, un tour de plus ajoute encore $a$ :
\[ x + na + a = x + (n+1)a . \]
\end{proof}
\source{https://github.com/Commutator-IO/learn-lean/blob/main/courses/01-college/07-algorithmique-et-programmation/AlgorithmiqueEtProgrammation.lean\#L50}{AlgorithmiqueEtProgrammation.lean\#L50}

\begin{theoreme}
Répéter $n$ fois le doublement multiplie par $2^n$ :
\[ \bigl(x \mapsto 2x\bigr)^{\,n}(x) = x \times 2^n . \]
\end{theoreme}

\begin{proof}
Même récurrence : la boucle vide laisse $x = x \times 2^0$, et un tour de plus multiplie
le résultat par deux, soit $x \times 2^n \times 2 = x \times 2^{n+1}$.
\end{proof}
\source{https://github.com/Commutator-IO/learn-lean/blob/main/courses/01-college/07-algorithmique-et-programmation/AlgorithmiqueEtProgrammation.lean\#L62}{AlgorithmiqueEtProgrammation.lean\#L62}

\subsection{Terminaison d'une boucle non bornée}

\begin{theoreme}
La boucle « retrancher $b$ tant que c'est possible » s'arrête, et ce qu'elle
laisse est le reste de la division :
\[ \mathrm{stq}(b, x) = x \bmod b \qquad (b > 0) . \]
\end{theoreme}

\begin{proof}
Par récurrence forte sur $x$, en distinguant les deux cas de la boucle.

Si $b \le x$, la boucle recommence sur $x - b$, strictement plus petit que $x$ puisque
$b > 0$ : l'hypothèse de récurrence s'applique et donne $(x - b) \bmod b$, qui vaut
$x \bmod b$ — retrancher le diviseur ne change pas le reste.

Sinon $x < b$, la boucle s'arrête et rend $x$, qui est bien le reste de la division de
$x$ par $b$.

C'est la décroissance stricte de $x$ à chaque tour qui fait la terminaison : une quantité
entière positive ne peut pas décroître indéfiniment.
\end{proof}
\source{https://github.com/Commutator-IO/learn-lean/blob/main/courses/01-college/07-algorithmique-et-programmation/AlgorithmiqueEtProgrammation.lean\#L78}{AlgorithmiqueEtProgrammation.lean\#L78}

\begin{figure}[h]
\centering
% Construction : droite graduée de 0 à 17, a = 17, b = 5.
%   Trois soustractions de b : 17 → 12 → 7 → 2. On s'arrête à 2 < 5, le reste.
%   Les arcs marquent les sauts de longueur b.
% SVG : x = 20 + 16x, y = 100 (axe).
\begin{tikzpicture}[scale=0.55, >=stealth, every node/.style={font=\small}]
  \draw[->] (-0.6, 0) -- (18.4, 0);
  \foreach \n in {0,5,10,15} { \draw (\n, 0.3) -- (\n, -0.3); \node[below] at (\n, -0.4) {$\n$}; }
  \foreach \d/\a in {17/12, 12/7, 7/2}
    \draw[->] (\d, 0.4) to[bend right=38] (\a, 0.4);
  \foreach \p in {17,12,7,2} \fill (\p, 0) circle (2.6pt);
  \node[above] at (14.5, 1.9) {$-b$};
  \node[above] at (9.5, 1.9) {$-b$};
  \node[above] at (4.5, 1.9) {$-b$};
  \node[below] at (2, -1.1) {reste};
  \node[below] at (17, -1.1) {$a$};
\end{tikzpicture}

\caption{Retrancher $b$ tant que c'est possible, ici $b = 5$ à partir de $a = 17$. Les sauts s'arrêtent dès qu'on passe sous $b$, et ce qui reste est le reste de la division euclidienne.}
\end{figure}

\begin{theoreme}
Pour $b = 0$, la condition d'entrée est fausse et la boucle rend $x$ sans rien
faire : $\mathrm{stq}(0, x) = x$.
\end{theoreme}

\begin{proof}
La condition « $0 < b$ et $b \le x$ » est fausse dès que $b = 0$ : la boucle n'entre
jamais. C'est cette hypothèse $b > 0$ qui garantit la décroissance dans le théorème
précédent.
\end{proof}
\source{https://github.com/Commutator-IO/learn-lean/blob/main/courses/01-college/07-algorithmique-et-programmation/AlgorithmiqueEtProgrammation.lean\#L93}{AlgorithmiqueEtProgrammation.lean\#L93}

\subsection{Correction d'un programme de calcul}

\begin{theoreme}
Correction du programme $A$ : il calcule le double de son entrée,
\[ \mathrm{ex}(A)(x) = 2x . \]
\end{theoreme}

\begin{proof}
On suit les trois instructions :
\[ x \ \longmapsto\ x + 3 \ \longmapsto\ 2(x+3) = 2x + 6 \ \longmapsto\ 2x + 6 - 6 = 2x . \]
\end{proof}
\source{https://github.com/Commutator-IO/learn-lean/blob/main/courses/01-college/07-algorithmique-et-programmation/AlgorithmiqueEtProgrammation.lean\#L100}{AlgorithmiqueEtProgrammation.lean\#L100}

\subsection{Deux programmes de longueurs différentes calculent la même chose}

\begin{theoreme}
Les deux programmes calculent la même chose :
$\mathrm{ex}(A)(x) = \mathrm{ex}(B)(x)$ pour toute entrée $x$.
\end{theoreme}

\begin{proof}
Le premier calcule $2x$ d'après le théorème précédent, le second aussi, en une seule
instruction.
\end{proof}
\source{https://github.com/Commutator-IO/learn-lean/blob/main/courses/01-college/07-algorithmique-et-programmation/AlgorithmiqueEtProgrammation.lean\#L108}{AlgorithmiqueEtProgrammation.lean\#L108}

\begin{theoreme}
Les deux programmes n'ont pourtant pas la même longueur : $A$ compte trois
instructions, $B$ une seule.
\end{theoreme}

\begin{proof}
Il suffit de compter les instructions : $3 \ne 1$. L'équivalence porte sur ce que les
programmes calculent, pas sur la façon dont ils le calculent.
\end{proof}
\source{https://github.com/Commutator-IO/learn-lean/blob/main/courses/01-college/07-algorithmique-et-programmation/AlgorithmiqueEtProgrammation.lean\#L115}{AlgorithmiqueEtProgrammation.lean\#L115}

\vfill
\noindent\rule{0.35\linewidth}{0.4pt}\par
\noindent{\footnotesize Leonardo de Moura et Sebastian Ullrich,
\emph{The Lean~4 Theorem Prover and Programming Language},
\emph{Automated Deduction — CADE~28}, Springer, 2021, p.~625--635,
\href{https://doi.org/10.1007/978-3-030-79876-5_37}{doi:10.1007/978-3-030-79876-5\_37}.\par}

\end{document}
